\chapter{Kết quả và đánh giá}

\section{Kết quả đạt được}

Sau quá trình phát triển và tích hợp, hệ thống ShopNexus đã hoàn thành
đầy đủ các mục tiêu đề ra. \Cref{tab:results} tổng hợp kết quả chính
của từng thành phần.

\begin{table}[htbp]
\centering
\caption{Tổng hợp kết quả các thành phần}
\label{tab:results}
\begin{tabular}{>{\raggedright\arraybackslash}p{0.25\linewidth} >{\raggedright\arraybackslash}p{0.65\linewidth}}
\toprule
\rowcolor{headerblue} \textcolor{white}{\textbf{Thành phần}} & \textcolor{white}{\textbf{Kết quả}} \\
\midrule
\rowcolor{rowgray} ESP32-CAM & Dual-core: audio 16 kHz liên tục + video JPEG streaming, điều khiển camera real-time (21 tham số) \\
Khử nhiễu DTLN & 1,8M tham số, 3,9 MB, 32 ms latency, huấn luyện trên VIVOS tiếng Việt \\
\rowcolor{rowgray} Voice Rephraze & STT qua gpt-4o-transcribe + sinh mô tả qua GPT-4o-mini, 4+ phong cách viết \\
Phân loại sản phẩm & Accuracy 90,14\%, F1 90,00\%, 32 danh mục, 45 ms/request (CPU) \\
\rowcolor{rowgray} Server trung gian & Relay TCP$\leftrightarrow$WebSocket, proxy API, spectrogram visualization \\
\bottomrule
\end{tabular}
\end{table}

\section{Kết quả thiết bị nhúng ESP32}

Thiết bị ESP32-CAM hoạt động ổn định trong môi trường thử nghiệm với kết nối
WiFi trong phạm vi 10--15 mét. Kiến trúc dual-core đảm bảo stream audio 16 kHz
không bị gián đoạn ngay cả khi camera đang chụp ở độ phân giải cao (UXGA
1600$\times$1200). Thời gian khởi động từ lúc cấp nguồn đến khi sẵn sàng
stream là khoảng 3--5 giây, bao gồm kết nối WiFi và khởi tạo camera.

Giao thức điều khiển camera cho phép thay đổi 21 tham số real-time từ giao
diện web, với thời gian phản hồi dưới 100 ms. Double buffering cho video
streaming đạt frame rate ổn định 10--15 FPS ở QVGA và 5--8 FPS ở VGA.

\section{Kết quả khử nhiễu DTLN}

Mô hình DTLN sau khi huấn luyện trên dữ liệu VIVOS kết hợp nhiễu Microsoft DNS
đạt khả năng khử nhiễu hiệu quả trên giọng nói tiếng Việt. Với chỉ 1,8 triệu
tham số và kích thước 3,9 MB, mô hình xử lý mỗi khung audio (32 ms) trong
thời gian tương đương, đảm bảo xử lý real-time.

\begin{figure}[htbp]
\centering
\fbox{\parbox{0.85\linewidth}{\centering\vspace{3cm}
\textit{Thêm hình: So sánh spectrogram trước và sau khử nhiễu --- hiển thị
2 spectrogram song song: (a) audio gốc có nhiễu nền (quạt, tiếng ồn đường phố),
(b) audio sau khi qua DTLN với nhiễu nền được loại bỏ rõ rệt,
giọng nói được giữ nguyên.}
\vspace{3cm}}}
\caption{So sánh spectrogram trước và sau khử nhiễu DTLN}
\label{fig:spectrogram-comparison}
\end{figure}

\section{Kết quả nhận dạng giọng nói và sinh mô tả}

Dịch vụ Voice Rephraze sử dụng gpt-4o-transcribe cho kết quả nhận dạng tiếng
Việt chính xác, kể cả với giọng vùng miền (Bắc, Trung, Nam). Module sinh mô
tả sản phẩm qua GPT-4o-mini tạo ra văn bản mạch lạc, tự nhiên, phù hợp với
phong cách viết được chọn.

\begin{figure}[htbp]
\centering
\fbox{\parbox{0.85\linewidth}{\centering\vspace{3cm}
\textit{Thêm hình: Ví dụ minh họa kết quả --- hiển thị: (1) text STT gốc
(lộn xộn, có từ lấp), (2) mô tả sản phẩm đã sinh theo phong cách ``Chuyên nghiệp'',
(3) mô tả theo phong cách ``Thân thiện''. Cho thấy sự khác biệt
giữa các phong cách viết.}
\vspace{3cm}}}
\caption{Ví dụ kết quả nhận dạng giọng nói và sinh mô tả sản phẩm}
\label{fig:stt-gen-example}
\end{figure}

\section{Kết quả phân loại sản phẩm}

Mô hình XLM-RoBERTa đạt accuracy 90,14\% và F1-score 90,00\% trên tập kiểm thử
32 danh mục. Thời gian suy luận trung bình 45 ms trên CPU cho single inference,
đủ nhanh cho ứng dụng real-time. \Cref{tab:classify-results} trình bày chi tiết
hiệu năng trên các cấu hình khác nhau.

\begin{table}[htbp]
\centering
\caption{Hiệu năng phân loại trên các cấu hình}
\label{tab:classify-results}
\begin{tabular}{>{\raggedright\arraybackslash}p{0.35\linewidth} >{\centering\arraybackslash}p{0.2\linewidth} >{\centering\arraybackslash}p{0.2\linewidth}}
\toprule
\rowcolor{headerblue} \textcolor{white}{\textbf{Cấu hình}} & \textcolor{white}{\textbf{Latency}} & \textcolor{white}{\textbf{Throughput}} \\
\midrule
\rowcolor{rowgray} Single inference (CPU) & $\approx$ 45 ms & 22 req/s \\
Batch 10 items (CPU) & $\approx$ 180 ms & 55 items/s \\
\rowcolor{rowgray} Single inference (GPU) & $\approx$ 15 ms & 66 req/s \\
Batch 100 items (GPU) & $\approx$ 800 ms & 125 items/s \\
\bottomrule
\end{tabular}
\end{table}

\begin{figure}[htbp]
\centering
\fbox{\parbox{0.85\linewidth}{\centering\vspace{2.5cm}
\textit{Thêm hình: Biểu đồ top-5 predictions cho một ví dụ cụ thể ---
bar chart hiển thị tên danh mục và confidence score (ví dụ:
``Electronics'' 87\%, ``Computers'' 6\%, ``Mobile Phones'' 3\%, ...).}
\vspace{2.5cm}}}
\caption{Ví dụ kết quả phân loại sản phẩm với top-5 predictions}
\label{fig:classify-example}
\end{figure}

\section{Kết quả tích hợp hệ thống}

Toàn bộ pipeline hoạt động liền mạch từ khi người dùng bắt đầu nói đến khi
kết quả hiển thị trên giao diện web. Thời gian xử lý end-to-end trung bình
khoảng 5--8 giây cho một đoạn mô tả 10--15 giây, bao gồm: truyền audio
từ ESP32 (real-time), khử nhiễu DTLN ($\approx$ 1 giây cho đoạn 10 giây),
STT ($\approx$ 2--3 giây), sinh mô tả ($\approx$ 1--2 giây), và phân loại
($\approx$ 45 ms).

\begin{figure}[htbp]
\centering
\fbox{\parbox{0.85\linewidth}{\centering\vspace{3cm}
\textit{Thêm hình: Ảnh chụp toàn bộ giao diện web đang hoạt động ---
hiển thị đồng thời: video stream từ camera, waveform âm thanh, spectrogram
trước/sau khử nhiễu, kết quả STT, mô tả sản phẩm đã sinh, và danh sách
top-K danh mục. Chú thích các vùng chức năng.}
\vspace{3cm}}}
\caption{Giao diện web hệ thống ShopNexus trong quá trình hoạt động}
\label{fig:full-demo}
\end{figure}

\section{Ưu điểm}

\begin{itemize}
\item
  \textbf{Pipeline end-to-end hoàn chỉnh:} Từ thu âm ESP32 $\rightarrow$ khử nhiễu
  $\rightarrow$ nhận dạng giọng nói $\rightarrow$ sinh mô tả $\rightarrow$ phân loại
  sản phẩm, hoạt động liền mạch không cần can thiệp thủ công.
\item
  \textbf{DTLN nhỏ gọn và hiệu quả:} Chỉ 3,9 MB (1,8M tham số), xử lý real-time
  32 ms, phù hợp triển khai trên thiết bị tài nguyên hạn chế. Huấn luyện trực
  tiếp trên dữ liệu tiếng Việt VIVOS đảm bảo hiệu quả với ngôn ngữ mục tiêu.
\item
  \textbf{Dual-domain denoising:} Kết hợp STFT (miền tần số) và Conv1D Encoder
  (miền đặc trưng học được), khử nhiễu toàn diện hơn so với phương pháp đơn miền.
  Giai đoạn 1 loại bỏ nhiễu tổng thể, giai đoạn 2 tinh chỉnh và phục hồi chi tiết.
\item
  \textbf{Phân loại chính xác:} XLM-RoBERTa đạt 90,14\% accuracy trên 32 danh mục,
  hỗ trợ đa ngôn ngữ nhờ nền tảng cross-lingual --- xử lý tốt cả mô tả tiếng Việt
  lẫn tên sản phẩm tiếng Anh.
\item
  \textbf{Dual-core ESP32 tối ưu:} Tách riêng audio/video trên 2 core đảm bảo
  không gián đoạn stream. 16 DMA buffer cho audio tạo vùng đệm 1024 ms,
  chống buffer underrun hiệu quả.
\item
  \textbf{Sinh mô tả đa phong cách:} Hỗ trợ 4+ phong cách viết (văn minh,
  chuyên nghiệp, thân thiện, văn hóa dân tộc), mở rộng được qua API.
  Prompt engineering xử lý tốt text lộn xộn, từ lấp từ STT.
\item
  \textbf{Kiến trúc microservices:} Mỗi dịch vụ phát triển, triển khai và
  scale độc lập. Có thể thay thế từng thành phần (ví dụ: đổi STT engine)
  mà không ảnh hưởng các thành phần khác.
\end{itemize}

\section{Hạn chế}

\begin{itemize}
\item
  \textbf{DTLN chỉ mono channel:} Chưa khai thác beamforming hay mảng microphone
  để tăng SNR đầu vào, hạn chế hiệu quả trong môi trường nhiều nguồn nhiễu.
\item
  \textbf{Phụ thuộc cloud cho STT:} Voice Rephraze sử dụng OpenAI API,
  yêu cầu kết nối internet ổn định và phát sinh chi phí vận hành. Không hoạt
  động được trong môi trường offline.
\item
  \textbf{DTLN nhạy cảm nhiễu phi tĩnh:} Hiệu năng giảm đáng kể với nhiễu
  biến đổi nhanh (tiếng nói chồng chéo, âm nhạc phức tạp) do LSTM cần thời
  gian thích nghi với pattern nhiễu mới.
\item
  \textbf{Pha không tối ưu:} Giai đoạn 1 của DTLN tái sử dụng pha gốc thay vì
  ước lượng pha sạch, có thể gây artifact âm thanh khi nhiễu nặng làm biến
  dạng pha ban đầu.
\item
  \textbf{ESP32 phụ thuộc WiFi:} Thiết bị không có cơ chế buffer offline ---
  nếu mất WiFi, dữ liệu audio/video bị mất hoàn toàn thay vì được lưu trữ
  tạm để gửi lại sau.
\item
  \textbf{Chưa có hệ thống đánh giá tự động:} Kết quả STT và sinh mô tả
  chưa được đánh giá bằng metrics chuẩn (WER, BLEU) do thiếu bộ test
  benchmark cho tiếng Việt trong ngữ cảnh thương mại điện tử.
\end{itemize}
